\subsection*{3. Differentiation Is a Counterfactual, Not a Feature Count}

The important question is not whether a buyer likes your product. It is whether, with your product available, their best decision becomes different.

Suppose an engineer is satisfied with a free tutorial and an internal reviewer. Your detailed curriculum does not move their choice. Another engineer has been asked to approve a production AI workflow and lacks a colleague who can challenge the architecture. A rigorous, bounded technical review may change that person's decision. The relevant difference is not necessarily curriculum coverage; it is the availability of competent critique at a consequential moment.

Define the buyer's real alternative before describing your superiority. A useful test is: \textbf{If our company vanished tomorrow, what would this person actually do?} The answer could be a direct competitor, an internal employee, a spreadsheet, a self-directed project, a cheaper substitution, continued suffering, or doing nothing. The founder who ignores the default alternative systematically exaggerates the value of their offering.

A differentiated product should be contrasted with the \emph{best feasible alternative} for the buyer, not the weakest rival you can find. In education, the strongest alternative may be doing actual work while consulting free technical material. In software, it may be modifying an existing workflow rather than purchasing new software. In a grocery store, it may be a trusted, good-enough brand already on the shelf.

\subsection*{4. A Model of the Buyer's Choice}

For conceptual purposes, write the perceived attractiveness of choice $j$ for buyer $i$ in circumstance $s$ as:

\begin{equation}
\begin{aligned}
U_{ijs} =\ &(B_{ijs} + E_{ijs} + I_{ijs} + A_{ijs}) \\
&-\ (P_{ijs} + T_{ijs} + R_{ijs} + F_{ijs}).
\end{aligned}
\end{equation}

Here, $B$ is expected functional or economic benefit, $E$ emotional benefit, $I$ identity or social value, and $A$ assurance or option value; $P$ is total monetary sacrifice, $T$ time and effort, $R$ perceived downside risk, and $F$ switching, learning, or approval friction. This is a \emph{reasoning device}, not a validated additive law: the terms overlap, interact, and are not generally measured in commensurable units.

The buyer chooses among accessible, permitted options, often without doing anything resembling an explicit calculation. For a high-stakes software purchase, formal scoring may approximate the process. For a familiar snack, habit and shelf accessibility may dominate. For a prestigious degree, social meaning and signal value may interact with academic learning and career options. One formula cannot flatten these differences.

Three conditions are frequently necessary: \textbf{access} (your offer is encountered), \textbf{preference} (the buyer likes its trade-off relative to alternatives), and \textbf{permission} (they have budget, authority, time, and acceptable risk). Preference without permission produces praise without purchases. Permission without preference produces procurement inertia. Access without relevance produces traffic without demand.

\subsection*{5. Different Buyers Want Different Trade-Offs}

A company cannot be best at everything because product attributes create resource claims and trade-offs. A working professional may prefer a demanding assessment even if it is inconvenient, because it produces trusted evidence. A casual learner may prefer flexibility over accountability. A corporate buyer may prefer transparent procurement terms over an exciting curriculum. A manufacturer may prefer predictable delivery to theoretically superior materials.

\textbf{The unit of differentiation is not the feature; it is the favorable trade-off for a recognizable customer situation.} A feature is useful only if it changes that trade-off. The first discipline is to ask what the buyer will willingly sacrifice to obtain the supposed advantage: extra money, setup time, admission effort, delivery speed, customization, privacy, a familiar brand, or flexibility.

A customer who says, ``I want all those benefits for free,'' has expressed a wish. A customer who selects your more demanding program over an easier, cheaper substitute has revealed a priority. Founders need to observe trade-offs, not merely collect adjectives such as \emph{innovative} or \emph{premium}.

\subsection*{6. Minimum Requirements, Choice Drivers, and Decoration}

Offer attributes generally fall into 3 operational roles for a particular decision:

\begin{itemize}
    \item \textbf{Table stakes} are the threshold qualities needed to be considered. A paid technical program needs intellectually competent teaching, honest terms, basic reliability, and a workable schedule. A B2B AI tool needs essential security and access controls. Failure here eliminates you before the comparison begins.
    \item \textbf{Choice drivers} are attributes where relative performance can actually change selection. These might be a rigorous review process, exceptionally low implementation effort, reliable local stock, an outstanding sensory experience, or credible proof of lower operating risk.
    \item \textbf{Decoration} includes attractive additions that do not materially affect purchase, use, or retention enough to justify their costs.
\end{itemize}

The same attribute can migrate between these roles as a category matures. A video-recorded lesson may once have delighted online learners; today it is simply assumed. Ask whether each proposed feature belongs to one of these categories for the intended buyer now, grounded in actual buying decisions rather than internal brainstorming.
